\documentclass[UTF8]{ctexart}

% 支持从项目根目录、强化学习目录或当前文档目录读取共享样式。
\makeatletter
\def\input@path{{../}{../../}}
\makeatother
\usepackage{researchnotes}

\title{马尔可夫决策过程与强化学习的关系}
\author{}
\date{}

\begin{document}
\maketitle

\section{问题模型：MDP 描述什么}
\label{sec:relation-mdp}

让机器人学会走迷宫，需要分清两个问题：迷宫任务的规则是什么，以及机器人怎样学会选择路线。
\keyterm{马尔可夫决策过程}（Markov decision process，MDP）是描述这类决策问题的数学模型；
\keyterm{强化学习}（reinforcement learning，RL）研究如何根据行动带来的反馈，学习更好的决策方式。
两者分别属于\keyterm{问题模型}与\keyterm{学习方法}这两个层次。

MDP 描述的基本过程是：决策者看到当前情况，采取动作，随后环境发生变化并给出奖励。
其中，决策者称为\keyterm{智能体}（agent），与它交互的外部系统称为\keyterm{环境}（environment）。
以地图固定、没有钥匙和电量限制的简单迷宫为例，需要说明：
\begin{itemize}
  \item \keyterm{状态}（state）：机器人当前所在的位置。
  \item \keyterm{动作}（action）：机器人可以选择的移动方向。
  \item \keyterm{转移概率}（transition probability）：选择某个方向后，到达各个位置的概率。
  \item \keyterm{奖励}（reward）：每次行动得到的数值反馈，例如普通移动得 $-1$ 分，进入出口得 $10$ 分并结束任务。
\end{itemize}
还应说明从哪里开始，以及怎样计算长期收益；例如把各步奖励相加，或让较远的奖励占较小权重。

“马尔可夫”要求：已知当前状态和动作后，过去的历史不再提供预测下一状态与奖励所需的额外信息。
在上述迷宫中，位置可以满足这一要求；若移动还取决于电量，则状态中还需包含剩余电量。
MDP 描述了任务的规则和收益，但\keyterm{不会自动给出好的行动方案}。

\section{学习方法：强化学习怎样利用经验}
\label{sec:relation-learning}

\keyterm{策略}（policy）规定智能体怎样选择动作。例如，一条策略可以规定在每个位置走哪个方向，
也可以规定以多大概率选择各个方向。强化学习的目标是根据经验改进策略，使长期累计奖励的期望更高，
而不只是让眼前这一步得分更高；绕开墙壁时，暂时远离出口也可能是合理选择。

机器人可以反复尝试，每次记录“原来的状态、采取的动作、得到的奖励、到达的状态”，
再利用这些经验改进行动选择。这里要区分：\keyterm{环境遵循某种规律，不等于智能体已经知道这些规律}。
任务可以由一个 MDP 描述，而机器人仍需通过交互了解行动的后果。

\keyterm{强化学习算法}（reinforcement learning algorithm）规定具体怎样利用经验。
例如，Q-learning 根据经验更新对动作长期收益的估计，再据此选择动作。
算法是学习步骤，策略是行动规则；二者都不是 MDP 本身。
同一个 MDP 可以采用不同算法学习，同一种算法也可以用于不同的 MDP。

\section{“实现 MDP”与“用强化学习求解”的区别}
\label{sec:relation-implementation}

如果把迷宫的移动和奖励规则写成可运行的模拟程序，可以称为\keyterm{实现一个 MDP 环境}。
如果让机器人在其中尝试，并根据反馈学习策略，则是在\keyterm{使用强化学习}。
前者把任务规则变成程序，后者利用经验学习如何行动。
因此，强化学习不是 MDP 的一种实现或特殊情况，强化学习算法也不是一种 MDP 模型。

更准确的表述是：\keyterm{强化学习常用于求解 MDP 所描述的决策问题，即通过经验学习好的策略}。
这里的“求解”是任务目标，不意味着任何算法都保证找到最优策略。
MDP 也可以用其他方法求解：若转移和奖励规律已知，可以直接利用模型进行\keyterm{规划}（planning）。
强化学习也研究只能观察部分环境信息等更一般的任务，因此两者的研究范围并不相同。

\end{document}
